\documentclass[11pt]{article}
\usepackage[utf8]{inputenc}
\usepackage{amsmath,amssymb}
\usepackage{authblk}
\usepackage[margin=1in]{geometry}
\usepackage[hidelinks]{hyperref}
\usepackage{xurl}
\usepackage{setspace}

\title{A Stationary-Action Principle for Real Quantum Mechanics,\\
and the Equation of Motion of the Free Boundary}
\author[1]{Claes Johnson}
\affil[1]{KTH Royal Institute of Technology, SE-100 44 Stockholm, Sweden \\
\texttt{claesjohnson@gmail.com}}
\date{\today}

\begin{document}
\maketitle

\begin{abstract}
Real Quantum Mechanics (RealQM) represents a system of $N$ electrons by $N$ charge densities carried on
\emph{non-overlapping} domains of ordinary three-dimensional space, the domains meeting at free boundaries
whose position is part of the solution. Its stationary states are obtained by minimising the Coulomb energy
over both the densities and the domains. We propose that the \emph{time-dependent} theory is governed by a
single \emph{stationary-action} principle in which the dynamical variables are the fields \emph{and} the
moving free boundaries, and we carry out the variation. Varying the fields returns the complex
time-dependent RealQM (Schr\"odinger) equations; varying the domains returns the law of motion of the free
boundary. Because the action is invariant under time translation and under an independent phase rotation of
each domain, N\"other's theorem guarantees, for free, both conservation of the total energy and conservation
of the charge \emph{of each electron separately}. The latter fixes the character of the interface: it is a
\emph{material} boundary across which no charge flows, so it moves with the normal charge-fluid velocity,
with that velocity continuous across it. This is the dynamical expression of RealQM's individuation of
electrons by the regions they occupy. The static minimiser is recovered as the stationary-state limit. We
give the boundary law explicitly, note that it corrects two naive prescriptions (Dirichlet with
slope-matching; Neumann with density-continuity) that violate one conservation law or the other, and
delimit where it is needed --- real-time, non-adiabatic multi-domain dynamics --- and where it is not
(Born--Oppenheimer chemistry; the timing of radioactive decay).
\end{abstract}

\setstretch{1.04}

\section{Introduction}
\label{sec:intro}

RealQM reformulates quantum mechanics for an $N$-electron system not as one complex amplitude on a
$3N$-dimensional configuration space but as $N$ real-space charge densities \cite{RealQM}. Each electron
$i$ is a density $\psi_i^2$ carried on its own region $\Omega_i\subset\mathbb R^3$, the regions
non-overlapping, $\Omega_i\cap\Omega_j=\varnothing$; the total charge density is $\rho=\sum_i q_i\psi_i^2$
with $q_i=-1$, and the ground state minimises the Coulomb energy over the fields and the domains, the
free boundaries $\partial\Omega_i$ being determined by the minimisation. This static variational
formulation reproduces atomic and molecular structure, and --- read for protons and electrons at
femtometre scale --- nuclear binding, all from electrostatics \cite{RealNucleus}.

Two things are then missing. First, a \emph{dynamical} principle: what equation governs the fields, and
--- the question we settle here --- what governs the \emph{motion of the free boundary} in time? Second, the
assurance that the dynamics respects the conservation laws it must: the energy, the total charge, and,
since the electrons are distinct objects individuated by their domains, the charge of \emph{each} electron.
We show that a single stationary-action principle supplies all of this at once. The free boundary is not an
auxiliary construction with an \emph{ad hoc} rule of motion; it is a dynamical variable of the same action
that gives the field equations, and its law of motion follows by variation, with the conservation laws
guaranteed by the symmetries of the action.

\section{The action}
\label{sec:action}

Work in atomic units ($\hbar=1$, $m_e=1$). The dynamical variables are the complex fields $\psi_i(x,t)$,
each supported on a time-dependent domain $\Omega_i(t)$, and the domains themselves (equivalently their
boundaries). Take the action
\begin{equation}
S=\int dt\;\Bigg\{\sum_i\int_{\Omega_i(t)}\!\Big[\tfrac{i}{2}\big(\psi_i^*\partial_t\psi_i-
\partial_t\psi_i^*\,\psi_i\big)-\tfrac12|\nabla\psi_i|^2\Big]dx-\tfrac12\!\iint\frac{\rho(x)\rho(y)}{|x-y|}
\,dx\,dy\Bigg\},
\label{eq:action}
\end{equation}
with $\rho=\sum_i q_i|\psi_i|^2$ and the constraints $\Omega_i\cap\Omega_j=\varnothing$ and, for each $i$,
$\int_{\Omega_i}|\psi_i|^2=1$ (unit charge per electron). The first bracket is the standard
Schr\"odinger Lagrangian density for each domain; the last term is the RealQM Coulomb energy, written as a
potential energy in the action. We take \eqref{eq:action} to be the basic principle of the theory: the
physical evolution is a stationary point of $S$ under variation of the fields and of the domains.

\section{Field variation: the RealQM equations}
\label{sec:field}

Vary $\psi_i^*$ at fixed domains. The bulk Euler--Lagrange equation is the complex time-dependent RealQM
(Schr\"odinger) equation
\begin{equation}
i\,\partial_t\psi_i=-\tfrac12\nabla^2\psi_i+q_i\,\phi[\rho]\,\psi_i,\qquad
\phi[\rho](x)=\int\frac{\rho(y)}{|x-y|}\,dy,
\label{eq:tdse}
\end{equation}
$\phi$ being the self-consistent Coulomb potential of the total density. The variation also leaves a term
on the free boundary, $-\int_{\partial\Omega_i}\!\operatorname{Re}(\partial_n\psi_i^*\,\delta\psi_i)$; its
treatment is dictated by the domain variation of the next section, to which it is naturally paired.

\section{Domain variation: the free boundary is a material interface}
\label{sec:boundary}

Now vary the domains. It is enough to consider two domains sharing an interface $\Gamma(t)$ with unit
normal $n$ pointing from $\Omega_1$ into $\Omega_2$; let the interface move with normal speed $V_n$. Two
facts about \eqref{eq:action}, both consequences of its symmetries, fix the interface law completely.

\paragraph{Per-electron charge conservation.} The action \eqref{eq:action} is invariant under an
\emph{independent} global phase rotation of each field, $\psi_i\to e^{i\alpha_i}\psi_i$ (the Coulomb term
depends only on the densities $|\psi_i|^2$). By N\"other's theorem each domain carries a conserved charge,
and the accompanying local law is the continuity equation
\begin{equation}
\partial_t|\psi_i|^2+\nabla\!\cdot\mathbf j_i=0,\qquad \mathbf j_i=\operatorname{Im}(\psi_i^*\nabla\psi_i).
\label{eq:cont}
\end{equation}
The charge of electron $i$ is $Q_i=\int_{\Omega_i(t)}|\psi_i|^2$, and with the boundary moving,
\begin{equation}
\frac{dQ_i}{dt}=\int_{\Omega_i}\partial_t|\psi_i|^2\,dx+\int_{\partial\Omega_i}|\psi_i|^2\,V_n\,dS
=\int_{\Gamma}\!\big(V_n\,|\psi_i|^2-\mathbf j_i\!\cdot\! n\big)\,dS ,
\end{equation}
using \eqref{eq:cont} and vanishing current on the outer boundary. For $Q_i$ to be conserved --- for each
electron to keep its unit charge, as its very identity in RealQM requires --- the integrand must vanish
pointwise on $\Gamma$:
\begin{equation}
\boxed{\;V_n=\frac{\mathbf j_i\!\cdot\! n}{|\psi_i|^2}=\mathbf v_i\!\cdot\! n\;}
\label{eq:material}
\end{equation}
where $\mathbf v_i=\mathbf j_i/|\psi_i|^2=\nabla S_i$ is the charge-fluid (Madelung) velocity of domain $i$,
with $\psi_i=|\psi_i|e^{iS_i}$. Equation \eqref{eq:material} says the interface is a \emph{material}
boundary: no charge crosses it, and it is carried along by the flow. Applying it to \emph{both} domains
forces the velocities to agree on the interface,
\begin{equation}
\mathbf v_1\!\cdot\! n=\mathbf v_2\!\cdot\! n=V_n\quad\text{on }\Gamma,
\label{eq:vmatch}
\end{equation}
so that a single boundary speed is well defined. Equations \eqref{eq:material}--\eqref{eq:vmatch} are the
law of motion of the free boundary: it advects with the (common) normal charge-fluid velocity.

\paragraph{The paired field condition.} The boundary term left by the field variation
(\S\ref{sec:field}) is now fixed consistently: on a material interface the natural condition is that the
fields carry no current \emph{through} the boundary in its own frame, which \eqref{eq:material} already
enforces; the residual variation vanishes with the density and its normal flux matched by
\eqref{eq:vmatch}. No charge is created, destroyed, or exchanged at $\Gamma$; the two clouds remain in
contact and simply move together.

\section{Conservation of energy and momentum}
\label{sec:noether}

Because \eqref{eq:action} carries no explicit time dependence, it is invariant under time translation, and
N\"other's theorem gives a conserved energy
\begin{equation}
E=\sum_i\int_{\Omega_i(t)}\tfrac12|\nabla\psi_i|^2\,dx+\tfrac12\!\iint\frac{\rho\rho}{|x-y|},
\label{eq:E}
\end{equation}
\emph{exactly}, for the coupled field--boundary evolution --- provided the boundary obeys its own
Euler--Lagrange law \eqref{eq:material} rather than an externally imposed rule. This is the crux: energy
conservation is not an extra property to be checked but a structural consequence of deriving the boundary
motion from the same action as the fields. Likewise, invariance under spatial translation gives conserved
total momentum $\mathbf P=\sum_i\int_{\Omega_i}\operatorname{Im}(\psi_i^*\nabla\psi_i)$. Thus one principle
delivers, together, the field equations, the boundary law, and the conservation of energy, momentum, and
per-electron charge.

\section{The static limit, and two naive prescriptions it corrects}
\label{sec:static}

For a stationary state $\psi_i(x,t)=\varphi_i(x)e^{-i\varepsilon_it}$ the currents $\mathbf j_i$ vanish, so
by \eqref{eq:material} $V_n=0$: the free boundary is at rest, and \eqref{eq:action} reduces to the static
Coulomb-energy minimisation over fields and domains that defines the RealQM ground state. The action
principle is therefore the dynamical parent of the static formulation used for atoms, molecules and nuclei
\cite{RealQM,RealNucleus}, not a separate construction.

The material-interface law also settles, and corrects, two prescriptions one is tempted to impose by hand.
A \emph{Dirichlet} interface ($\psi_i=0$ on $\Gamma$) with the boundary placed where the normal slopes
match, $|\partial_n\psi_1|=|\partial_n\psi_2|$, conserves per-electron charge but is \emph{not} the
stationary-action law and does not, in general, conserve energy. A \emph{Neumann} interface
($\partial_n\psi_i=0$) placed where the densities are continuous, $|\psi_1|=|\psi_2|$, can be arranged to
conserve energy but transfers charge across the moving boundary and so \emph{violates} per-electron charge
conservation. Only the material law \eqref{eq:material}--\eqref{eq:vmatch}, in which the boundary moves with
the flow, respects both conservation laws, precisely because both descend from the one action.

\section{Scope}
\label{sec:scope}

The boundary law matters exactly where the domains genuinely rearrange \emph{in time}. In
\emph{Born--Oppenheimer} chemistry the electrons re-minimise at each nuclear geometry, using the static
free boundary of \S\ref{sec:static}; the time-dependent law is not invoked, and ordinary reaction paths,
barriers and binding energies are unaffected. It is needed for \emph{real-time, non-adiabatic}
multi-domain dynamics --- attosecond electron motion, charge migration, photochemistry, collisions --- where
how the electron territories reshape in time is itself the physics. It is, by contrast, immaterial to
phenomena that are insensitive to the partition: the timing and rate of radioactive decay, for instance,
are carried by the total charge density tunnelling a barrier and do not depend on the interior boundary
between domains at all.

\section{Conclusion}
\label{sec:concl}

RealQM is governed by a single stationary-action principle, \eqref{eq:action}, whose dynamical variables are
the charge-density fields and the free boundaries between them. Varying the fields gives the complex
time-dependent RealQM equations; varying the domains gives the law of motion of the free boundary, which the
independent phase symmetry of each domain fixes to be a \emph{material} interface, advected by the common
normal charge-fluid velocity, \eqref{eq:material}--\eqref{eq:vmatch}. Energy, momentum and per-electron
charge are then conserved automatically, by N\"other's theorem, because all of it descends from one
symmetric action. The static energy minimisation that defines atoms, molecules and nuclei is the
stationary-state limit of the same principle. The free boundary, in short, is not an appendage to RealQM
but one of its dynamical variables, and it moves --- like everything else in the theory --- by stationary
action.

\begin{thebibliography}{9}
\bibitem{RealQM} C.~Johnson, \emph{Quantum Mechanics as Multiphase 3D Continuum Mechanics}, manuscript,
2026; \url{https://claes542.github.io/RealMolecule/gallery.html}.
\bibitem{RealNucleus} C.~Johnson, \emph{RealNucleus: Nuclear Binding as Dual Coulomb Confinement, without a
Strong or Weak Force}, manuscript, 2026.
\bibitem{ComplexRealQM} C.~Johnson, \emph{Complex Time-Dependent RealQM: Currents, Radiation, and Magnetism
from the Real-Space Continuum}, manuscript, 2026.
\end{thebibliography}

\end{document}
